\SourceLeaf{246}%
also im Werk von \textsc{Lobatschewski} nicht gefunden, doch die Entwicklung
folgt einem anderen Weg als dem, den ich selbst eingeschlagen habe, und zwar führt sie
\textsc{Lobatschewski} meisterhaft und in echt geometrischem Geist aus.
Ich glaube, Sie auf das Buch aufmerksam machen zu müssen, das Ihnen gewiss
ganz besonderen Genuss bereiten wird. \mbox{.\ .\ .\ .\ .\ .}

\par\bigskip\Needspace{18\baselineskip}

\begin{center}
\textsc{Gauss} an W. \textsc{Struve}.\quad Göttingen, 11. Dezember 1846.
\end{center}

\mbox{.\ .\ .\ .\ .\ .\ .} Ebenso danke ich Ihnen sehr für die übrigen Sendungen.
Für die russischen Arbeiten von
\textsc{Lobatschewski} bin ich wohl in erster Linie Ihrem Sohn dankbar, dem gegenüber
ich vor einigen Jahren, als er hier war, meinen Wunsch ausgesprochen hatte.
Bitte empfehlen Sie mich seinem freundlichen Andenken. Meine
Russischkenntnisse sind wohl etwas eingerostet,
da ich seit mehr als einem Jahr nicht dazu gekommen bin, auch nur
einen russischen Buchstaben anzusehen. Ich hoffe jedoch, bei der ersten freien
Gelegenheit das Versäumte rasch nachzuholen, und werde dann der Lektüre jener
interessanten Aufsätze besondere Aufmerksamkeit widmen. Die kleine
deutsche Schrift von \textsc{Lobatschewski} besaß ich schon vorher
selbst. \mbox{.\ .\ .\ .\ .\ .\ .}


